% 分野別類題: 定数係数2階線形非同次微分方程式（未定係数法） 解答例
% コンパイル: TEXINPUTS="../../../00_common:$TEXINPUTS" latexmk -xelatex answers.tex
\documentclass[a4paper,11pt]{article}
\usepackage{preamble}
\usepackage{shoakubox}

\begin{document}

\kamokumei{類題演習 解答例 --- 定数係数2階線形非同次微分方程式（未定係数法）}

\daimon{【解法】未定係数法による特殊解の求め方と導出}

\noindent
定数係数2階線形非同次微分方程式
\[
y'' + a\,y' + b\,y = f(x)
\]
の一般解は、対応する同次方程式 $y'' + a y' + b y = 0$ の一般解（同次解）$y_h$ と、非同次方程式の特殊解 $y_p$ の和
\[
y = y_h + y_p
\]
で与えられる。これは、$y_1, y_2$ がともに非同次方程式の解ならば $y_1 - y_2$ は同次方程式の解になることから従う（線形微分作用素 $L[y] = y'' + ay' + by$ の線形性 $L[y_h+y_p] = L[y_h]+L[y_p] = 0 + f(x) = f(x)$ による）。

\medskip
\noindent
\textbf{同次解}は特性方程式 $r^2 + ar + b = 0$ の根 $r_1, r_2$ から、
\begin{itemize}
\item $r_1 \neq r_2$（実根）: $y_h = C_1 e^{r_1 x} + C_2 e^{r_2 x}$
\item $r_1 = r_2 = r$（重根）: $y_h = (C_1 + C_2 x)e^{rx}$
\item $r = \alpha \pm i\beta$（複素根）: $y_h = e^{\alpha x}(C_1\cos\beta x + C_2 \sin\beta x)$
\end{itemize}
と定まる。

\medskip
\noindent
\textbf{特殊解}は、$f(x)$ の形に応じて次のように仮定形を置き、方程式に代入して未定係数を決定する（未定係数法）。
\begin{itemize}
\item $f(x)$ が $n$ 次多項式 $\Rightarrow$ $y_p = $（$n$ 次多項式）
\item $f(x) = k e^{\mu x}$ $\Rightarrow$ $y_p = A e^{\mu x}$
\item $f(x) = \sin\omega x$ や $\cos\omega x$ $\Rightarrow$ $y_p = A\cos\omega x + B\sin\omega x$
\end{itemize}
ただし、\textbf{仮定した特殊解の形が同次解 $y_h$ の項と重なる場合（共鳴）}は、そのままでは代入すると恒等的に $0 = f(x)$ の形になり係数が定まらない。この場合は、仮定形に $x$ を掛けて同次解と一次独立にする。すなわち、仮定形が同次解に含まれる重複度を $m$ とすると、仮定形全体に $x^{m}$ を掛ける（例えば重根 $r$ に対して $f(x)=ke^{rx}$ のときは $y_p = Ax^{2}e^{rx}$ とする）。

\medskip
\noindent
係数決定後は、必ず元の方程式に代入して検算する。初期条件が与えられている場合は、一般解 $y = y_h + y_p$（および必要なら $y'$）に代入して積分定数を定める。

\clearpage
\begin{mondaiboxS}{問1}
$y'' - y = x^{2} + 1$ の一般解を求めよ。
\end{mondaiboxS}
\kaitou
\textbf{同次解:} 特性方程式 $r^{2} - 1 = 0$ より $r = \pm 1$。よって
\[
y_h = C_{1}e^{x} + C_{2}e^{-x}
\]

\noindent
\textbf{特殊解:} $f(x) = x^{2}+1$ は2次多項式で、同次解に多項式項は含まれないので共鳴なし。$y_p = Ax^{2}+Bx+C$ とおくと $y_p'' = 2A$ より
\[
y_p'' - y_p = 2A - (Ax^{2}+Bx+C) = -Ax^{2} - Bx + (2A - C)
\]
これが $x^{2}+1$ に等しいので
\[
-A = 1,\qquad -B = 0,\qquad 2A - C = 1
\]
よって $A=-1,\ B=0,\ C = 2A-1 = -3$。すなわち $y_p = -x^{2}-3$。

（検算: $y_p''=-2$、$y_p''-y_p = -2-(-x^2-3) = x^2+1$ OK）

\noindent
したがって一般解は
\[
\boxed{\; y = C_{1}e^{x} + C_{2}e^{-x} - x^{2} - 3 \;}
\qquad (C_1, C_2 \text{ は任意定数})
\]

\clearpage
\begin{mondaiboxS}{問2}
$y'' + 4y = \sin 2x$ の一般解を求めよ。
\end{mondaiboxS}
\kaitou
\textbf{同次解:} 特性方程式 $r^{2}+4=0$ より $r = \pm 2i$。よって
\[
y_h = C_{1}\cos 2x + C_{2}\sin 2x
\]

\noindent
\textbf{特殊解（共鳴あり）:} 素朴には $y_p = A\cos 2x + B\sin 2x$ とおきたいが、これは同次解と一次従属（共鳴）であるため、$x$ 倍した
\[
y_p = x(A\cos 2x + B\sin 2x)
\]
を仮定する。
\[
y_p' = A\cos 2x - 2Ax\sin 2x + B\sin 2x + 2Bx\cos 2x
\]
\[
y_p'' = -4A\sin 2x - 4Ax\cos 2x + 4B\cos 2x - 4Bx\sin 2x
\]
したがって
\[
y_p'' + 4y_p = \bigl(-4A\sin 2x - 4Ax\cos 2x + 4B\cos 2x - 4Bx\sin 2x\bigr) + 4x(A\cos 2x + B\sin 2x)
= -4A\sin 2x + 4B\cos 2x
\]
（$x\cos 2x,\ x\sin 2x$ の項は打ち消し合う。）これが $\sin 2x$ に等しいので
\[
-4A = 1,\qquad 4B = 0 \quad\Longrightarrow\quad A = -\frac{1}{4},\ B = 0
\]
よって $y_p = -\dfrac{x}{4}\cos 2x$。

（検算: 上式に $A=-1/4,\,B=0$ を代入すると $y_p''+4y_p = -4A\sin2x+4B\cos2x = \sin 2x$ OK）

\noindent
したがって一般解は
\[
\boxed{\; y = C_{1}\cos 2x + C_{2}\sin 2x - \frac{x}{4}\cos 2x \;}
\qquad (C_1, C_2 \text{ は任意定数})
\]

\clearpage
\begin{mondaiboxS}{問3}
$y'' - 2y' + y = e^{x}, \quad y(0) = 1,\ y'(0) = 0$ を解け。
\end{mondaiboxS}
\kaitou
\textbf{同次解:} 特性方程式 $r^{2}-2r+1 = (r-1)^{2} = 0$ より重根 $r=1$。よって
\[
y_h = (C_{1}+C_{2}x)e^{x}
\]

\noindent
\textbf{特殊解（重根との共鳴）:} $f(x)=e^{x}$ は $e^{x}$ 型であるが、$e^{x}$ は同次解に（重複度2で）含まれるため、単純な $Ae^x$ でも $Axe^x$ でも共鳴してしまう。そこで
\[
y_p = Ax^{2}e^{x}
\]
とおく。
\[
y_p' = A e^{x}(x^{2}+2x), \qquad y_p'' = A e^{x}(x^{2}+4x+2)
\]
したがって
\[
y_p'' - 2y_p' + y_p = Ae^{x}\bigl[(x^{2}+4x+2) - 2(x^{2}+2x) + x^{2}\bigr] = Ae^{x}\cdot 2 = 2Ae^{x}
\]
これが $e^{x}$ に等しいので $2A=1$、すなわち $A = \dfrac{1}{2}$。よって $y_p = \dfrac{1}{2}x^{2}e^{x}$。

\noindent
\textbf{一般解と初期条件:}
\[
y = (C_{1}+C_{2}x)e^{x} + \frac{1}{2}x^{2}e^{x} = e^{x}\left(C_{1} + C_{2}x + \frac{1}{2}x^{2}\right)
\]
\[
y' = e^{x}\left(C_{1} + C_{2}x + \frac{1}{2}x^{2}\right) + e^{x}(C_{2}+x)
\]
$y(0) = C_{1} = 1$。$y'(0) = C_{1}+C_{2} = 0$ より $C_{2} = -1$。

\noindent
したがって解は
\[
\boxed{\; y = \left(1 - x + \frac{1}{2}x^{2}\right)e^{x} \;}
\]

（検算: $y' = e^x\left(\dfrac{1}{2}x^2\right)$、$y'' = e^x\left(\dfrac{1}{2}x^2+x\right)$ より
$y''-2y'+y = e^x\left[\left(\dfrac{1}{2}x^2+x\right) - 2\cdot\dfrac{1}{2}x^2 + \left(1-x+\dfrac{1}{2}x^2\right)\right] = e^x\cdot 1 = e^x$ OK。
また $y(0)=1$, $y'(0)=\dfrac{1}{2}\cdot 0^2=0$ OK）

\end{document}
